% ============================================================
%  Poul Martin Møller,
%  "Om Begrebet Ironie"
%
%  A fragment (Brudstykke II) occasioned by Provost Tryde's
%  review of Sibbern's Æsthetik (Maanedsskrift for Litteratur,
%  bd. 13, 1835). Møller, as a member of the journal's editorial
%  board, added aphoristic remarks on the concept of irony but
%  "bragte kun Indledningen til Ende" (see the editor's note):
%  only the treatment of MORAL irony was completed; the piece
%  breaks off as it turns to poetic irony.
%
%  Published posthumously in:
%  Efterladte Skrifter af Poul M. Møller, bd. 3 (Tredie Bind),
%  ed. F.C. Olsen. Kjøbenhavn: C.A. Reitzel.
%  (Transcribed from the 3rd edition, "Tredie Udgave".)
%  Printed pp. 152--158. PDF (bibliotek scan) pp. 163--169;
%  offset PDF = printed + 11.
%
%  TRANSCRIPTION (Danish)
%  Conventions: original orthography kept (aa-spellings,
%  capitalised nouns, æ/ø, "philosophisk", "Conseqvents",
%  "Qvantum", the -us endings Egoismus/Nihilismus/Idealismus as
%  printed). Fraktur long-s rendered as s. Danish quotes „…"
%  kept. Book titles -> \emph{}. Printed page numbers -> \opage{N}
%  (marginal, at the point the printed page begins).
% ============================================================

\documentclass[12pt]{article}
\usepackage[utf8]{inputenc}
\usepackage[T1]{fontenc}
\usepackage{graphicx}
% Old Danish "id est" mark (reversed c + colon), printed as "ɔ:"
\DeclareUnicodeCharacter{0254}{\reflectbox{c}}
\usepackage{libertinus}
\usepackage{libertinust1math}
\usepackage{textalpha}
\usepackage[danish]{babel}
\usepackage[protrusion=true,expansion=true]{microtype}
\usepackage[margin=1.4in]{geometry}
\usepackage{setspace}
\usepackage{fancyhdr}
\usepackage{marginnote}
\usepackage[colorlinks=true,linkcolor=black,urlcolor=black,
  pdftitle={Om Begrebet Ironie},
  pdfauthor={Poul Martin Møller}]{hyperref}
\setlength{\headheight}{15pt}
\pagestyle{fancy}
\fancyhf{}
\fancyhead[LE,RO]{\thepage}
\fancyhead[RE]{\textit{Om Begrebet Ironie}}
\fancyhead[LO]{\textit{Poul Martin Møller}}
\setlength{\parindent}{1.5em}
\setlength{\parskip}{0pt}
\onehalfspacing

\newcommand{\opage}[1]{\marginnote{\footnotesize\textit{[#1]}}}
\newcommand{\ombreak}{\par\medskip\begin{center}\rule{0.32\textwidth}{0.4pt}\end{center}\medskip}

\title{\textbf{Om Begrebet Ironie}}
\author{Poul Martin Møller}
\date{Efterladte Skrifter, bd.\ 3 (Brudstykke II)}

\begin{document}
\maketitle

\begin{center}\rule{0.32\textwidth}{0.4pt}\end{center}
\medskip

\opage{152}I Anledning af de Oplysninger, som den ærede Forfatter af forestaaende Anmeldelse\footnote{Der menes Provst Trydes Anmeldelse af Sibberns Æsthetik i Maanedsskrift f.\ Lit.\ 13 Bd.\ (1835) og navnlig p.\ 200 flg. Ved som Medlem af Maanedsskriftets Redaction at gjøre sig bekjendt med Provst Trydes Manuskript, fik P.\ M.\ Lyst til at vedføie sine Bemærkninger; men bragte kun Indledningen til Ende. Da den kan betragtes som et Hele for sig (om den moralske Ironie), turde den fortjene en Plads i hans Skrifter, saameget mere som den afgiver et Bidrag til det omhandlede Begrebs Historie i vor Litteratur; dog maatte det erindres, at Gjenstandens Behandlingsmaade har været betinget af Fordringerne til en Indledning.} har meddeelt om Begrebet Ironie, saaledes som det har dannet sig i Nutidens æsthetiske Talebrug, være det et af Redactionens Medlemmer tilladt at tilføie nogle aphoristiske Bemærkninger, der maaskee kunne give endeel Læsere en tydeligere Indsigt i dette Begrebs Oprindelse og forskjellige Forvandlinger.

Det undskyldes, at Begyndelsen hentes noget langt borte; vi skulle desuagtet snart igjen komme tilbage til Sagen. I Grækernes moralske Selvbevidsthed saaledes som vi lære den at kjende i dens første Spire hos Folkets ældste Digtere, f.\ Ex.\ af flere Steder hos Homer, og saaledes som vi siden forefinde den i fuldstændig Udvikling hos Aristoteles, er den naturlige Begjering betragtet som det Godes positive Side; Fornuften derimod er det begrændsende Princip, som kun billiger et vist Qvantum af dette \opage{153}naturlige Indhold og forkaster et andet. Dyden er et Maadehold, en Middelvei mellem to Yderligheder. Vi forefinde altsaa paa dette Standpunct to forskjellige Bestemmelser i det Godes Begreb, der forliges med hinanden, fordi deres Modsætning ikke ret tydelig er blevet udviklet. Men det maatte nødvendigviis derefter bringes til Bevidsthed, at den blotte Drift som en blind, ubegribelig og ufornuftig Bestemmelsesgrund ikke kunde billiges af Fornuften. Dette maa allerede Aristoteles, den fuldkomneste Repræsentant for hiint Standpunct, have følt, og deraf maa man vel forklare sig den Synderlighed, at den practiske Philosophies Indhold efter hans Begreber om Erkjendelse ikke kan blive Gjenstand for Videnskab i Ordets høiere Betydning.

Af de to heterogene og formeentlig uforsonlige Elementer i det Godes Begreb har nu vexelviis snart det ene, snart det andet med Eensidighed været forfægtet i Videnskabens Historie. Snart har man aldeles opgivet Forsøget at bestemme Livets Øiemed ved fornuftig Tænkning og anseet det for Viisdom at lade den naturlige Begjering i sin tilfældige Skikkelse være den ledende Magt, snart har man ved en Tænkning, der holdt sig blot paa det individuelle Standpunct, villet udlede de forefundne moralske Maximer af en dogmatisk Sætning eller af en aldeles abstract og indholdsløs Regel. Saaledes har man seet de aldeles modsigende Forskrifter: „Lader os vende tilbage til Naturstanden" (Rousseau), og: „Man maa forlade den naturlige Forfatning" (Hobbes) --- at blive gjorte gjeldende efter hinanden.

Den Bestræbelse, at bestemme det Godes Indhold blot fra den rene Fornufts Standpunct eller ved en blot apriorisk Tænkning, har været gjennemført med størst Conseqvents af den britiske Philosophies Hovedmænd. I denne Philosophie har man forsøgt af blot formelle Principer at udlede den i Samfundet gjeldende Kreds af sædelige Maximer. Al Heteronomie, enhver blot naturlig Bestemmelsesgrund vilde man forkaste og ved reen \opage{154}apriorisk Tænkning specificere den moralske Selvbevidstheds Indhold. Men uagtet den Begeistring, hvormed Tanken om absolut Autonomie i en saadan Betydning fyldte denne Skoles Tilhængere, kunde dog en saadan Plan umuligen blive videnskabeligen gjennemført. Baade Kant og Fichte see sig aabenbar nødsagede til at betjene sig af heteronomiske Maximer for at skaffe deres abstracte Principer et virkeligt Indhold.

Idet man nu lærte, at de naturlige Formaal skulde udføres som Midler for Moralloven, men ikke for den Lyst, der uundgaaeligen forbandt sig med deres Udførelse, maatte der ifølge Sagens Natur søges en Methode, efter hvilken det lod sig afgjøre, hvilket Formaal der i ethvert enkelt Moment var det rette Middel, hvorved den almindelige Lov kunde sættes i Kraft. Her lagde nu Fichte den høieste Fuldmagt i Individets Bevidsthed, saa at dette efter egen moralsk Overbeviisning i det enkelte Tilfælde skulde afgjøre, hvad der var dets Pligt. Dette Standpunct, der altid maa beholde sin (relative) Gyldighed, tillægger altsaa Individet Evne til umiddelbart at erkjende, efter hvilken af de forskjellige Regler det Almindelige i det enkelte Tilfælde skal realiseres. ---

Derfor kan der af denne moralske Idealisme, naar den fastholdes paa en eensidig Maade og ikke forenes med Erkjendelse af, at Fornuftbud allerede realiseres udenfor Individet, drages forkastelige Conseqventser. Den subjective Overbeviisning bliver da betragtet som det Høieste, idet Individets Villie identificeres med Moralloven. Den individuelle Villie stilles da over enhver speciel Lov og enhver bestemt Regels døde Bogstav, der kun har Gyldighed, naar den sanctioneres af den individuelle Villie. Denne over Alt ophøiede autonomiske Villie betragtes da følgerigtigen som hævet over de borgerlige Love, der kuns, hvor de stemme med den, have Gyldighed, --- en Synsmaade, som tillader alle Forbrydelser, der kunne begaaes og ofte have været be\opage{155}gaaede med god Samvittighed og ungdommelig, moralsk Begeistring. Saaledes gives der Intet, som i sig selv og uafhængigt af det moralske Subject er godt, men dette, som identisk med Moralloven, kan hellige hvilkensomhelst Handling ved sin Overbeviisning; thi det, der ellers i sig selv skulde være slet, ophører at være det, naar Villien, som er Hovedsagen, er ædel, eller naar Hensigten var god.

Fra denne Theorie gjøres let Overgangen til en anden, der sætter den subjective Villie over Moralloven; thi den er jo egentlig practisk sat derover, naar det, som i sig selv er Slet, kan helliges af Villien. Moralloven har i Grunden tabt sin Hellighed og Individet gjør om ikke sædeligt, saa dog logisk Fremskridt ved udtrykkelig at antage sig hævet over Moralloven. Denne Theorie findes udtrykkelig udtalt i F.\ Schlegels ligesaa frække som geniale Ungdomsskrift \emph{Lucinde}, som er et af de frodigste Vandskud af den hensynsfrie, dristige Selvtillid', der er sig selv nok. Vi citere til Bekræftelse paa det Anførte et Par Yttringer af \emph{Lucinde} selv: „Intet er galere, end naar Moralister nu gjøre hinanden Bebreidelser for Egoismus! thi hvilken Gud kan være ærværdig for det Menneske, der ikke er sin egen Gud?" Fremdeles at „man maa i Tilværelsens Nydelse hæve sig over alle endelige Formaal og Forsætter"; at det Høieste altsaa er den guddommelige Dovenskab, og at man kuns i den ægte Passivitets hellige Slethed kan erindre sig sit hele Jeg, hvilken rene Vegeteren da er det egentlige Liv og den sande Religion. --- Dette Standpunct har Hegel optaget i sin practiske Philosophie og betegnet med Navnet Ironie, hvilket han forklarer som den „Subjectivitet, der veed sig selv som det Høieste." Heri erkjender han intet Positivt uden den Virksomhed, at udtale denne Bevidsthed for Andre, og det er jo ligefrem en Opgivelse af alle Bestræbelser i den objective Verden.

Uagtet denne practiske Resignation er uendelig forskjellig fra \opage{156}den energiske Moral, der læres i Fichtes Værker (hvilket det vel er overflødigt at bemærke), er det dog vist, at den nedstammer fra den fichteske Idealismus, som det ogsaa er bekjendt, at Fred.\ Schlegel var en af den fichteske Videnskabslæres meest enthusiastiske Beundrere. I Særdeleshed er Ironien en conseqvent Fortsættelse af den frugtesløse Stræben, at bringe et i sig afsluttet Moralsystem i Stand blot fra Individets Standpunct. Denne Fremgangsmaade maatte nødvendigviis ende med Savn af alt Indhold, med en moralsk Nihilismus.

Den her omtalte Ironie betegner altsaa ei et blot tænkeligt Standpunct i den moralske Bevidstheds Udvikling; det er heller ikke fremtraadt som et uforklarligt Indfald, hvorpaa en enkelt særsindet Mand tilfældigviis var bragt, men det var et nødvendigt Resultat af en vildfarende Tankegang, der viste sig under lignende Skikkelse hos flere Samtidige.

Selv hos os har denne i Tidens Udvikling dybt begrundede Forvildelse havt sin svagere Gjenklang, skjøndt den ifølge Folkets Beskedenhed og Forsigtighed kun har efterladt faa eller ingen Spor i vor Literatur. Men i Conversation har man oftere kunnet høre hiin Synsmaade klinge saa temmelig efter, og om den ikke har været udtalt med fuld Forvisning og Selvtillid, har den dog været yttret som et Resultat, man maatte føres til, naar man vilde overlade sig til en fordomsfri Tænknings Førelse. Den har meget almindelig viist sig som et Slags Mistroe til Begrebet Moralitet, hvilket man helst undgik i det videnskabelige Foredrag, da man ikke meer deri fandt et adæqvat Udtryk for den practiske Frihed; man har ogsaa engang imellem i Poesien spøgt med Begreberne Moral, Moralitet og moralsk, og yttret en dunkel Bevidsthed om, at det ikke hang rigtig sammen dermed. Ja En og Anden har ogsaa betjent sig af den dristige Vending, at sand Dygtighed ikke lod sig indskrænke af den snævre Morallov.

En Skikkelse af den moralske Bevidsthed, som er nær be\opage{157}slægtet med Ironien, er den sygelige Sentimentalitet: den kan i det Mindste føre til det samme Resultat, nemlig til den Mening, at Mennesket kan være hævet over sine jordiske Pligter. En ørkesløs Længsel efter et bedre Liv bringer da den Sentimentale til at opgive alle Livets bestemte Formaal, fordi Alt hernede kun er forgjængeligt og ikke afgiver nogen værdig Skueplads for den, der har Sands for noget Høiere. Der gjøres paa dette Standpunct en Adskillelse mellem et blot Endeligt paa den ene Side, som er aldeles udeelagtigt i Uendelighed, og et derfra ved et svælgende Dyb adskilt, aldeles ubestemt høiere Liv, som har været kaldt det Bedre eller det Evige. Da nu hiint er den Sentimentales eneste værdige Øiemed, mener han at maatte i passiv Resignation, med Opgivelse af al bestemt Bestræbelse, oppebie det Høieres Komme til ham. Man kan her anvende, hvad Fichte, saavidt erindres, etsteds har sagt: Mange holde for, at man blot behøver at lade sig begrave for at opnaae et saligt Liv. --- Denne Sentimentalitet ligner deri Ironien, at begge ansee sig for hævede over alle Pligtbud: denne, fordi han sætter sin individuelle Vilkaarlighed som det Høieste, hiin, fordi han betragter al endelig Handling i Tid og Rum som en lavere Virkekreds og troer, at han besmitter sit indre Livs Reenhed ved at give sig af med dens Berøring. Begge holde sig altsaa for gode til at have Pligter: Ironikeren, fordi hans individuelle Jeg er det Absolute; den Sentimentale, fordi hans Virkekreds i Livet ei er udenfor ham; han lever kuns et indre Liv, som fyldes med Længsler og Forudfølelse af det Tilkommende, som er det eneste Reelle. Men de ere deri forskjellige fra hinanden, at Ironikeren betragter sin Subjectivitet som det Høieste, den Sentimentale tager det Absolute som Et uden for sin Person, som udenfor den hele synlige Verden. Dog udjævnes denne Forskjel noget derved, at den Sentimentales Bevidsthed er Sædet for Erkjendelse af det Endeliges Intethed, af det ubestemte Høieres Realitet, hvorved han \opage{158}kommer til at staae langt over Alle, som med Iver forfølge bestemte practiske Formaal.

Ogsaa kan Ironikeren efter sit System, om han vil, uden at hindres af noget moralsk Snørliv, overlade sig til den groveste Sandseligheds Nydelse --- en Conseqvents, der erkjendes i \emph{Lucinde}. Dette strider egentlig imod den sentimentale Theorie, men da Læren her umulig kan følges i sin hele Conseqvents, og Handling umulig, saalænge Livet varer, kan aldeles opgives, og følgelig Berøring med Endeligheden ei kan aldeles undgaaes, kan den Sentimentale lettere end den, der hylder en fornuftig Sædelære, forsvare for sig selv, at han foretrækker det meest sandselige Slags af de i sig selv ligegyldige Handlinger. (Naar det ikke er Umagen værdt at gjøre det Gode, er det heller ikke Umagen værdt at lade det Slette.)

Men dette var et Sidespring, der blot skulde tjene til at tydeliggjøre den hidtil omtalte practiske Ironie ved at sammenbringe den med den anden Art af practisk Nihilisme, hvormed den i flere Henseender stemmer overeens. En af de mange Betydninger, hvori man har taget Ordet Ironie i den nyere Philosophie, synes hermed tilstrækkelig at være forklaret. Saaledes bruger Hegel Ordet Ironie og betegner altsaa dermed en af „det Ondes moralske Former." Hvad enten man nu anseer Ordet for vel eller ilde valgt, fortjener dets Gjenstand Videnskabens Opmærksomhed, og Pladsen, som den indtager i Bevidsthedens Historie, er den uimodsigelig rigtig anviist af Hegel. Vi gaae nu over til den poetiske Ironie. --- ---

\ombreak

\end{document}
